\scanpage{182}\mspage{37}
les ``complétés profinis'' \inserted{(au sens d'Artin-Mazur)} premiers à $p$ (la car. de $k$) de tels \unsure{espaces}. On voit clairement comment aborder un tel programme en car. $0$, mais on prévoit de l'amusement supplémentaire, voire du mystère, en car $p>0$, pour les variétés qui, même birationnellement, refuseraient de se laisser remonter en car. $0$!

\begin{center}\rule{0.18\linewidth}{0.4pt}\end{center}

De ces cogitations essentiellement géométriques je ne saurais à l'heure actuelle tirer un programme précis de développement de structures algébriques adéquates pour les exprimer. Je me bornerai à quelques remarques en marge.

\deleted{\unsure{a)|$\alpha$)}} Depuis longtemps j'ai été \unsure{intrigu-|intrigué} par l'idée d'une ``linéarisation'' d'un type d'homotopie (ordinaire), i.e. la question du type: Si $X$ est un type d'homotopie, combien d'information cohomologique du type: cohomologie de $X$ à coefficients variables $M$ (\unsure{constant|constants} ou constant tordus), structure multiplicative $\deleted{\unsure{+|=}}H^i(X,M)\times H^j(X,N)\unsure{-|\to}H^{i\unsure{\cdot|+}j}(X,M\otimes N)$, plus éventuellement d'autres opérations cohomologiques -- est-il nécessaire d'avoir pour réconstituer le type d'homotopie entièrement \deleted{\unsure{?}}? (disons,
% archive sheet number: 182
